% V009 only: B=(0,0), C=(a,0), A=(u,h), a=BC,b=CA,c=AB.
% Arguments: mode,scale,a,b,c,caption. Exact geometry, equal x/y scales.
% 0 opposite-edge coefficients; 1 equal-height areas; 2 two bisectors;
% 3 common-base height ratio; 4 coordinate-distance check; 5 numeric areas; 6 two divisions.
\providecommand{\VIncenterFigure}[6]{%
\begin{center}
\begin{tikzpicture}[scale=#2,font=\small,line cap=round,line join=round,>=Stealth]
\pgfmathsetmacro{\va}{#3}\pgfmathsetmacro{\vb}{#4}\pgfmathsetmacro{\vc}{#5}
\pgfmathsetmacro{\vL}{\va+\vb+\vc}
\pgfmathsetmacro{\vu}{(\vc*\vc+\va*\va-\vb*\vb)/(2*\va)}
\pgfmathsetmacro{\vh}{sqrt(\vc*\vc-\vu*\vu)}
\pgfmathsetmacro{\vix}{\va*(\vu+\vc)/\vL}
\pgfmathsetmacro{\vr}{\va*\vh/\vL}
\pgfmathsetmacro{\vd}{\va*\vc/(\vb+\vc)}
\pgfmathsetmacro{\vsa}{\va*\vr/2}
\pgfmathsetmacro{\vsb}{\vb*\vr/2}
\pgfmathsetmacro{\vsc}{\vc*\vr/2}
\coordinate (A) at (\vu,\vh);\coordinate (B) at (0,0);
\coordinate (C) at (\va,0);\coordinate (I) at (\vix,\vr);
\coordinate (D) at (\vd,0);
\coordinate (X) at ($(B)!(I)!(C)$);
\coordinate (Y) at ($(C)!(I)!(A)$);
\coordinate (Z) at ($(A)!(I)!(B)$);
\ifnum#1=1
 \fill[blue!14] (I)--(B)--(C)--cycle;
 \fill[orange!18] (I)--(C)--(A)--cycle;
 \fill[violet!13] (I)--(A)--(B)--cycle;
\fi
\ifnum#1=5
 \fill[blue!14] (I)--(B)--(C)--cycle;
 \fill[orange!18] (I)--(C)--(A)--cycle;
 \fill[violet!13] (I)--(A)--(B)--cycle;
\fi
\ifnum#1=2
 \fill[blue!9] (A)--(B)--(D)--cycle;
\fi
\ifnum#1=3
 \fill[blue!13] (I)--(B)--(C)--cycle;
\fi
\draw[thick] (A)--(B)--(C)--cycle;
\ifnum#1=0
 \draw[blue!70!black,very thick] (B)--node[midway,below=6pt] {$a$} (C);
 \draw[orange!80!black,very thick] (C)--node[midway,right=6pt] {$b$} (A);
 \draw[violet!80!black,very thick] (A)--node[midway,left=6pt] {$c$} (B);
 \draw[->,blue!70!black,thick] (I)--(A);
 \draw[->,orange!80!black,thick] (I)--(B);
 \draw[->,violet!80!black,thick] (I)--(C);
 \node[above=5pt,blue!70!black] at (A) {$A$};
 \node[below left=4pt,orange!80!black] at (B) {$B$};
 \node[below right=4pt,violet!80!black] at (C) {$C$};
\else
 \ifnum#1=4
  \draw[->,gray!75] (-1.5,0)--({\va+1.8},0) node[right] {$x$};
  \draw[->,gray!75] (0,-1.2)--(0,{\vh+1.5}) node[above] {$y$};
  \node[above=5pt] at (A) {$A(u,h)$};
  \node[below left=6pt] at (B) {$B(0,0)$};
  \node[below right=6pt] at (C) {$C(a,0)$};
 \else
  \node[above=5pt] at (A) {$A$};
  \node[below left=4pt] at (B) {$B$};
  \node[below right=4pt] at (C) {$C$};
 \fi
 \ifnum#1=2
  \path (A)--node[midway,left=6pt] {$c$} (B);
  \path (C)--node[midway,right=6pt] {$b$} (A);
 \else
  \ifnum#1=5
   \path (B)--node[midway,below=7pt] {$a=#3$} (C);
   \path (C)--node[pos=.65,right=5pt] {$b=#4$} (A);
   \path (A)--node[pos=.35,left=5pt] {$c=#5$} (B);
  \else
   \ifnum#1=6\else
    \path (B)--node[pos=.85,below=7pt] {$a$} (C);
   \fi
   \path (C)--node[pos=.7,right=5pt] {$b$} (A);
   \path (A)--node[pos=.3,left=5pt] {$c$} (B);
  \fi
 \fi
\fi
\ifnum#1=1
 \foreach \T in {A,B,C}{\draw[dashed,gray!65] (I)--(\T);}
 \draw[blue!60!black] (I) circle[radius=\vr];
 \foreach \F/\T in {X/C,Y/A,Z/B}{%
  \draw[blue!75!black,thick] (I)--(\F);
  \pic[draw,angle radius=2.7mm] {right angle=I--\F--\T};}
 \path (I)--node[midway,right=4pt] {$r$} (X);
 \path (I)--node[midway,above left=2pt] {$r$} (Z);
 \path (I)--node[pos=.4,above left=4pt] {$r$} (Y);
\fi
\ifnum#1=5
 \foreach \T in {A,B,C}{\draw[dashed,gray!65] (I)--(\T);}
 \foreach \F/\T in {X/C,Y/A,Z/B}{%
  \draw[blue!75!black] (I)--(\F);
  \pic[draw,angle radius=2.7mm] {right angle=I--\F--\T};}
 \node at ({\va/4},{\vr/3}) {$\pgfmathprintnumber[fixed,precision=0]{\vsa}$};
 \node at ({(\vix+2*\va+\vu)/4},{(\vr+\vh)/4}) {$\pgfmathprintnumber[fixed,precision=0]{\vsb}$};
 \node at ({(\vix+\vu)/4},{(\vr+\vh)/4}) {$\pgfmathprintnumber[fixed,precision=0]{\vsc}$};
\fi
\ifnum#1=2
 \draw[->,orange!80!black,thick] (I)--(A);
 \draw[->,blue!70!black,thick] (I)--(D);
 \draw[thick,violet!65!black] (B)--(I);
 \draw[blue!65!black,very thick] (B)--(D);
 \draw[orange!75!black,very thick] (D)--(C);
 \node[below=6pt] at (D) {$D$};
 \path (B)--node[midway,below=6pt] {$BD$} (D);
 \path (D)--node[midway,below=6pt] {$DC$} (C);
 \pic[draw,violet!65!black,angle radius=10mm] {angle=I--B--A};
 \pic[draw,violet!65!black,angle radius=10mm] {angle=D--B--I};
 \pic[draw,orange!75!black,angle radius=9mm] {angle=B--A--D};
 \pic[draw,orange!75!black,angle radius=9mm] {angle=D--A--C};
\fi
\ifnum#1=3
 \coordinate (H) at (\vu,0);\coordinate (K) at (\vix,0);
 \draw[orange!80!black,thick] (A)--(D);
 \draw[dashed,orange!80!black,thick] (A)--(H);
 \draw[blue!75!black,thick] (I)--(K);
 \pic[draw,angle radius=2.5mm] {right angle=A--H--D};
 \pic[draw,angle radius=2.5mm] {right angle=I--K--C};
 \node[below left=5pt] at (H) {$H$};
 \node[below=18pt] at (K) {$K$};
 \node[below right=5pt] at (D) {$D$};
 \fill (D) circle[radius=1.6pt];
 \path (A)--node[pos=.55,left=4pt,orange!80!black] {$AH$} (H);
 \path (I)--node[midway,right=12pt,blue!75!black] {$IK$} (K);
\fi
\ifnum#1=6
 \draw[orange!80!black,very thick] (A)--(I);
 \draw[blue!70!black,very thick] (I)--(D);
 \draw[blue!70!black,very thick] (B)--(D);
 \draw[orange!80!black,very thick] (D)--(C);
 \fill (D) circle[radius=1.6pt];\node[below=6pt] at (D) {$D$};
 \path (B)--node[midway,below=6pt] {$BD=13/2$} (D);
 \path (D)--node[midway,below=6pt] {$DC=15/2$} (C);
 \node[left=6pt] at ($(A)!.45!(I)$) {$AI$};
 \node[right=6pt] at ($(I)!.5!(D)$) {$ID$};
\fi
\ifnum#1=4
 \foreach \F/\T in {X/C,Y/A,Z/B}{%
  \draw[blue!75!black,thick] (I)--(\F);
  \pic[draw,angle radius=2.7mm] {right angle=I--\F--\T};}
 \path (I)--node[midway,right=9pt] {$d_{BC}$} (X);
 \path (I)--node[pos=.35,above=7pt] {$d_{AB}$} (Z);
 \path (I)--node[pos=.4,above=8pt] {$d_{AC}$} (Y);
 \draw[dashed,gray!65] (A)--(D);
 \fill (D) circle[radius=1.5pt];\node[below=18pt] at (D) {$D$};
\fi
\fill (I) circle[radius=2pt];
\ifnum#1=4
 \node[right=16pt] at (I) {$P$};
\else
 \ifnum#1=2
  \node[left=6pt] at (I) {$I$};
 \else
  \ifnum#1=3
   \node[above right=5pt] at (I) {$I$};
  \else
   \node[right=16pt] at (I) {$I$};
  \fi
 \fi
\fi
\end{tikzpicture}
\par\smallskip{\small #6}
\end{center}}

% V009 examples: each branch uses the exact vertices of that exercise.
\providecommand{\VIncenterExample}[1]{%
\begin{center}
\begin{tikzpicture}[font=\small,>=Stealth,line cap=round,line join=round]
\ifnum#1=1
 \begin{scope}[scale=1.35]
 \coordinate (A) at (0,3);\coordinate (B) at (4,0);\coordinate (C) at (0,0);\coordinate (I) at (1,1);
 \coordinate (X) at (1,0);\coordinate (Y) at (0,1);
 \draw[->,gray!75] (-.6,0)--(4.8,0) node[right] {$x$};
 \draw[->,gray!75] (0,-.5)--(0,3.8) node[above] {$y$};
 \fill[blue!7] (A)--(B)--(C)--cycle;\draw[thick] (A)--(B)--(C)--cycle;
 \draw[blue!70!black,thick] (I) circle[radius=1];
 \draw[blue!70!black,dashed] (I)--(1,0) (I)--(0,1);
 \pic[draw,angle radius=2.5mm] {right angle=I--X--B};
 \pic[draw,angle radius=2.5mm] {right angle=I--Y--A};
 \node[above left=3pt] at (A) {$A(0,3)$};\node[below=5pt] at (B) {$B(4,0)$};
 \node[below left=4pt] at (C) {$C(0,0)$};
 \path (A)--node[midway,above right=4pt] {$c=5$} (B);
 \path (C)--node[midway,below=7pt] {$a=4$} (B);
 \path (A)--node[midway,left=7pt] {$b=3$} (C);
 \fill[blue!70!black] (I) circle[radius=2pt];\node[above right=5pt] at (I) {$I$};
 \end{scope}
\fi
\ifnum#1=2
 \begin{scope}[scale=1.3]
 \coordinate (A) at (0,0);\coordinate (B) at (3,0);\coordinate (C) at (0,4);\coordinate (P) at (1,1);
 \draw[thick] (A)--(B)--(C)--cycle;
 \draw[blue!70!black] (P) circle[radius=1];
 \draw[->,blue!70!black,thick] (P)--(A);
 \draw[->,orange!80!black,thick] (P)--(B);
 \draw[->,violet!80!black,thick] (P)--(C);
 \node[below left=4pt,blue!70!black] at (A) {$A$};
 \node[below right=4pt,orange!80!black] at (B) {$B$};
 \node[above left=4pt,violet!80!black] at (C) {$C$};
 \draw[very thick,blue!70!black] (B)--node[midway,right=5pt] {$a=5$} (C);
 \draw[very thick,orange!80!black] (C)--node[midway,left=6pt] {$b=4$} (A);
 \draw[very thick,violet!80!black] (A)--node[midway,below=6pt] {$c=3$} (B);
 \fill (P) circle[radius=2pt];\node[right=7pt] at (P) {$P$};
 \end{scope}
\fi
\ifnum#1=3
 \begin{scope}[scale=1.4]
 \coordinate (A) at (0,0);\coordinate (B) at (4,0);\coordinate (C) at (0,3);\coordinate (D) at (12/7,12/7);
 \draw[->,gray!75] (-.6,0)--(4.8,0) node[right] {$x$};
 \draw[->,gray!75] (0,-.6)--(0,3.8) node[above] {$y$};
 \draw[thick] (A)--(B)--(C)--cycle;
 \draw[very thick,blue!70!black] (B)--node[midway,above right=5pt] {$BD$} (D);
 \draw[very thick,orange!80!black] (D)--node[midway,above right=5pt] {$DC$} (C);
 \draw[->,green!45!black,thick] (A)--(D);
 \node[below left=3pt] at (A) {$A(0,0)$};\node[below=5pt] at (B) {$B(4,0)$};
 \node[above left=3pt] at (C) {$C(0,3)$};
 \fill (D) circle[radius=2pt];\node[above right=8pt] at (D) {$D(12/7,12/7)$};
 \pic[draw,angle radius=8mm] {angle=B--A--D};\pic[draw,angle radius=8mm] {angle=D--A--C};
 \end{scope}
\fi
\ifnum#1=4
 \begin{scope}[scale=1.4]
 \coordinate (A) at (0,0);\coordinate (B) at (4,0);\coordinate (C) at (0,3);\coordinate (I) at (1,1);
 \draw[->,gray!75] (-.6,0)--(4.8,0) node[right] {$x$};
 \draw[->,gray!75] (0,-.6)--(0,3.8) node[above] {$y$};
 \draw[thick] (A)--(B)--(C)--cycle;
 \draw[blue!70!black] (I) circle[radius=1];
 \draw[->,green!45!black,thick] (A)--(3.3,3.3) node[above right=4pt] {$y=x$};
 \draw[->,blue!70!black,very thick] (A)--(1,0) node[below=6pt] {$\vec u$};
 \draw[->,orange!80!black,very thick] (A)--(0,1) node[left=6pt] {$\vec v$};
 \fill[blue!70!black] (I) circle[radius=2pt];\node[below right=5pt] at (I) {$I$};
 \fill[gray!85] (2.5,2.5) circle[radius=2pt];\node[right=5pt] at (2.5,2.5) {$Q$};
 \node[below left=3pt] at (A) {$A$};\node[below=5pt] at (B) {$B(4,0)$};
 \node[above left=3pt] at (C) {$C(0,3)$};
 \end{scope}
\fi
\ifnum#1=5
 \begin{scope}[scale=.8]
 \coordinate (A) at (0,0);\coordinate (B) at (6,0);\coordinate (C) at (0,8);
 \coordinate (I) at (2,2);\coordinate (G) at (2,8/3);\coordinate (M) at (3,4);
 \draw[->,gray!75] (-.7,0)--(7,0) node[right] {$x$};
 \draw[->,gray!75] (0,-.7)--(0,8.8) node[above] {$y$};
 \draw[thick] (A)--(B)--(C)--cycle;
 \draw[blue!70!black] (I) circle[radius=2];
 \draw[blue!70!black,dashed] (A)--(24/7,24/7);
 \draw[orange!85!black,dashed] (A)--(M) (C)--(3,0);
 \fill[blue!70!black] (I) circle[radius=2pt];\node[below=8pt,blue!70!black] at (I) {$I$};
 \fill[orange!85!black] (G) circle[radius=2pt];\node[left=10pt,orange!85!black] at (G) {$G$};
 \node[below left=3pt] at (A) {$A(0,0)$};\node[below=5pt] at (B) {$B(6,0)$};
 \node[above left=3pt] at (C) {$C(0,8)$};
 \fill (M) circle[radius=1.5pt];\node[right=4pt] at (M) {$M$};
 \end{scope}
\fi
\end{tikzpicture}
\par\smallskip{\small
\ifcase#1\or 内切圆与两坐标轴相切，$I=(r,r)=(1,1)$；第三边也必须与圆相切。
\or 三边与顶点同色配对：$10:8:6=2(5:4:3)$；约去整体倍数后，定位同一内心 $P=I$。
\or $D$ 从 $B$ 向 $C$ 走全程的 $4/7$；$BD:DC=4:3$，不是 $3:4$。
\or $I(1,1)$ 与 $Q$ 都在 $y=x$ 上；方向相同，位置可以不同。
\or 蓝色角平分线定位 $I(2,2)$，橙色中线相交于 $G(2,8/3)$；$M$ 是 $BC$ 的中点。
\fi}
\end{center}}

\begin{statementbox}
\textbf{一、这条结论解决什么问题？}

已知三角形三个顶点的坐标，怎样\textbf{不求两条角平分线的交点}，就直接求内心坐标？
若题目给出一个带有边长系数的向量等式，怎样\textbf{认出其中的点就是内心}？

\medskip
\textbf{二、使用条件与字母约定}

在真正的三角形 $ABC$ 中，设
\[
BC=a>0,\qquad CA=b>0,\qquad AB=c>0,
\]
并且任意两边之和大于第三边。$I$ 是\textbf{三条内角平分线的交点}，即三角形的内心。
始终注意：\textbf{$a$ 是 $A$ 的对边，$b$ 是 $B$ 的对边，$c$ 是 $C$ 的对边。}

\VIncenterFigure{0}{.49}{14}{15}{13}{同色配对：蓝边 $BC=a$ 对应顶点 $A$，橙边 $CA=b$ 对应 $B$，紫边 $AB=c$ 对应 $C$。}

\medskip
\textbf{三、核心结论（先理解第一式，其余都是变形）}
\begin{enumerate}
\item \textbf{内心向量等式：}
\[
\boxed{a\vec{IA}+b\vec{IB}+c\vec{IC}=\vec 0.}
\]
若点 $P$ 满足 $a\vec{PA}+b\vec{PB}+c\vec{PC}=\vec 0$，则 $P$ 就是内心 $I$（反向也成立）。
\item \textbf{任选参考点 $O$，直接定位内心：}
\[
\boxed{\vec{OI}=\frac{a\vec{OA}+b\vec{OB}+c\vec{OC}}{a+b+c}.}
\]
通俗地说：把三个顶点的位置分别乘以\textbf{它的对边长}，相加后除以三边总长。
\item \textbf{内心坐标：} 若 $A(x_1,y_1),B(x_2,y_2),C(x_3,y_3)$，则
\[
\boxed{I\left(\frac{ax_1+bx_2+cx_3}{a+b+c},\quad
\frac{ay_1+by_2+cy_3}{a+b+c}\right).}
\]
\item \textbf{角平分线与对边交点：} 若 $AD$ 平分 $\angle BAC$ 且 $D\in BC$，则
\[
\frac{BD}{DC}=\frac cb,\qquad
\boxed{\vec{BD}=\frac c{b+c}\vec{BC}},\qquad
\boxed{\vec{AD}=\frac b{b+c}\vec{AB}+\frac c{b+c}\vec{AC}}.
\]
\VIncenterFigure{2}{.48}{14}{15}{13}{$A,I,D$ 依次排列；$BD:DC=c:b$，且 $AI:ID=(b+c):a$。}
\item \begin{minipage}[t]{\linewidth}\textbf{只判断方向：} 由 $A$ 出发的\textbf{内角平分线}方向，与
\[
\boxed{\frac{\vec{AB}}{|\vec{AB}|}+\frac{\vec{AC}}{|\vec{AC}|}}
\]
同向。将中间的加号改为减号，得到\textbf{外角平分线所在直线}的方向。\textbf{方向向量只决定一条直线的方向，不能单独给出内心坐标。}
% Equal unit vectors u=(1,0),v=(1/2,sqrt(3)/2), drawn with one common scale.
\begin{center}
\begin{tikzpicture}[scale=4,font=\small,>=Stealth,line cap=round,line join=round]
\coordinate (A) at (0,0);\coordinate (U) at (1,0);
\coordinate (V) at (.5,{sqrt(3)/2});\coordinate (W) at (1.5,{sqrt(3)/2});
\coordinate (Z) at (.5,{-sqrt(3)/2});\coordinate (K) at (.75,{sqrt(3)/4});
\fill[blue!7] (A)--(U)--(W)--(V)--cycle;
\draw[dashed,gray!70] (U)--(W)--(V);
\draw[->,blue!75!black,thick] (A)--node[midway,below=5pt] {$\vec u$} (U);
\draw[->,orange!80!black,thick] (A)--node[midway,left=5pt] {$\vec v$} (V);
\draw[->,green!50!black,thick] (A)--node[pos=.8,above=6pt] {$\vec u+\vec v$} (W);
\draw[gray!75,dashed] (V)--(U);
\draw[->,violet!75!black,thick] (A)--node[midway,right=7pt] {$\vec u-\vec v$} (Z);
\pic[draw,angle radius=3mm] {right angle=A--K--U};
\pic[draw,green!50!black,angle radius=9mm] {angle=U--A--W};
\pic[draw,green!50!black,angle radius=9mm] {angle=W--A--V};
\node[left=5pt] at (A) {$A$};
\node[below right=4pt] at (U) {$U$};\node[above left=4pt] at (V) {$V$};
\node[above=4pt] at (W) {$W$};
\node[below=5pt] at (Z) {外角平分线方向};
\end{tikzpicture}
\par\smallskip{\small $|\vec u|=|\vec v|=1$，四边形 $AUWV$ 是菱形；两条对角线垂直，和向量平分内角。}
\end{center}

\end{minipage}
\end{enumerate}

\textbf{四、可先掌握的几何关系}

设 $r$ 是内切圆半径，$S$ 为 $\triangle ABC$ 的面积。连接 $IA,IB,IC$，则
\VIncenterFigure{1}{.48}{14}{15}{13}{三块面积分别对应底边 $a,b,c$；蓝色垂线都是半径 $r$，不是到顶点的距离。}
\[
S_{IBC}=\frac12 ar,\quad S_{ICA}=\frac12 br,\quad S_{IAB}=\frac12 cr,
\qquad \boxed{S_{IBC}:S_{ICA}:S_{IAB}=a:b:c.}
\]
\end{statementbox}
